\clearpage
\pdfbookmark[0]{Supplement to the Theory of the Combination of Observations}{supplementum}
\thispagestyle{empty}
\SourceLeaf{060}
\begin{center}
\vspace*{2em}
{\Large SUPPLEMENT\par}
\bigskip
{\Large TO THE THEORY\par}
\bigskip
{\Large\bfseries OF THE COMBINATION OF OBSERVATIONS\par}
\bigskip
{\large SUBJECT TO THE SMALLEST ERRORS\par}
\bigskip
BY\par
\bigskip
{\large CARL FRIEDRICH GAUSS\par}
\bigskip
{\small PRESENTED TO THE ROYAL SOCIETY OF SCIENCES IN 1826, ON SEPTEMBER 16.\par}
\bigskip
\rule{.7\linewidth}{.3pt}\par
\bigskip
Commentationes societatis regiae scientiarum Gottingensis recentiores.
\mbox{Vol.~\textsc{vi}.}
\par\medskip
Göttingen MDCCCXXVIII.
\par\bigskip
\rule{.7\linewidth}{.3pt}
\end{center}
\clearpage
\renewcommand{\Article}[1]{\par\medskip\Needspace{4\baselineskip}\pdfbookmark[1]{Art. #1}{supplementum-art-#1}\centerline{\textbf{#1.}}\nobreak\smallskip}
\SourceLeaf{062}
\begin{center}
{\large SUPPLEMENT\par}
\medskip TO THE THEORY\par
\medskip {\large OF THE COMBINATION OF OBSERVATIONS\par}
\medskip SUBJECT TO THE SMALLEST ERRORS.\par
\medskip\rule{.25\linewidth}{.3pt}
\end{center}
\Article{1}
In the treatment of the theory of the combination of observations included in Volume V of the Commentationes Recentiores, we assumed that the quantities whose values are given by observations lacking absolute precision depend on certain unknown elements in such a way that they are represented as given functions of these elements, and that the main task is to derive these elements from the observations as accurately as possible.

In most cases, indeed, this assumption applies immediately. In other cases, however, the problem presents itself in a slightly different form, so that at first sight it seems doubtful how it can be reduced to the required form. It often happens, namely, that the quantities to which the observations refer have not yet been represented as functions of certain elements, and do not even seem reducible to such a form, at least not conveniently or without circuitous procedures; while, on the other hand, the nature of the matter supplies certain conditions which the true values of the observed quantities must satisfy exactly and necessarily.

Nevertheless, on closer consideration, it is easily seen that this case does not in fact differ essentially from the other, but can be reduced to it. For if the number of observed quantities is denoted by \(\pi\) and the number of conditional equations by \(\sigma\), and \(\pi-\sigma\) of the former are chosen arbitrarily, nothing prevents us from taking these themselves as the elements and considering the others, whose num\SourceLeaf{063}%
ber will be \(\sigma\), as functions of them by means of the conditional equations; in this way the matter will have been reduced to our assumption.

Yet even if this route leads quite conveniently to the intended goal in very many cases, it cannot be denied that it is less natural, and that it is therefore worthwhile to treat the problem separately in this other form, all the more so because it admits a very elegant solution. Indeed, since this new solution leads to quicker calculations than the solution of the problem in its earlier form whenever \(\sigma\) is smaller than \(\frac12\pi\), or, what is the same, whenever the number of elements, denoted in the preceding paper by \(\rho\), is greater than \(\frac12\pi\), it will be appropriate in such a case to prefer the new solution, which we shall explain in the present paper, to the earlier one, provided that the conditional equations can be drawn from the nature of the problem without circuitous procedures.

\Article{2}
Let us denote by \(v,v',v''\) etc. the quantities, \(\pi\) in number, whose values become known through observation, and let an unknown quantity depend on them in such a way that it is represented by a given function of them, say \(u\); further, let \(l,l',l''\) etc. be the values of the derivatives
\[
\frac{\mathrm du}{\mathrm dv},\quad \frac{\mathrm du}{\mathrm dv'},\quad \frac{\mathrm du}{\mathrm dv''}\ \text{etc.}
\]
corresponding to the true values of the quantities \(v,v',v''\) etc. Just as substituting these true values in the function \(u\) gives its true value, so, if for \(v.\ v',v''\) etc. values differing from the true ones by the errors \(e,e',e''\) etc. respectively are substituted, an erroneous value of the unknown quantity will be obtained, whose error can be written
\[
=le+l'e'+l''e''+\text{etc.}
\]
provided, as we shall always assume, that the errors \(e,e',e''\) etc. are so small that (when the function \(u\) is not linear) their squares and products may be neglected. And although the sizes of the errors \(e,e',e''\) etc. remain uncertain, the uncertainty inherent in such a determination of the unknown quantity can nevertheless be estimated in general, namely by the mean error to be feared in such a determination, which by the principles of the preceding paper becomes
\[
=\BareRoot{(llmm+l'l'm'm'+l''l''m''m''+\text{etc.})}
\]
\SourceLeaf{064}%
where \(m,m',m''\) etc. denote the mean errors of the observations, or, if each observation is subject to equal uncertainty,
\[
=m\BareRoot{(ll+l'l'+l''l''+\text{etc.})}
\]
It is clear that in this calculation it is equally legitimate to take for \(l,l',l''\) etc. the values of the derivatives corresponding to the observed values of the quantities \(v,v',v''\) etc.

\clearpage
